\exercisesection{1}

¶ 1. (a) 证明群 G 不能是两个异于 G 的子群的并。证明对每个至少含两个元素的集合 I，交换群 \(G = \mathbf{F}_{2}^{(I)}\) 是三个异于 G 的子群的并。

\begin{enumerate}
\def\labelenumi{(\alph{enumi})}
\setcounter{enumi}{1}
\item
  设 \((\mathbf{H}_i)_{i\in I}\) 是群 \(G\) 的子群的一个有限族，使得每个 \(\mathbf{H}_i\) 都是 \(G\) 的无限指数子群。证明 \(G\) 不能是诸 \(\mathbf{H}_i\) 的有限多个左陪集的并。（对 \(I\) 中元素的个数用归纳法论证；若存在两个不同的指标 \(i,j\) 使指数 \((\mathbf{H}_i:(\mathbf{H}_i\cap \mathbf{H}_j))\) 有限，则可删去 \(\mathbf{H}_i\) ；反之，若 \(\mathbf{H}_i\cap \mathbf{H}_j\) 在 \(\mathbf{H}_i\) 中都有无限指数（对每对不同指标 \((i,j)\) ），则取一个指标 \(k\) ，使 \(\mathbf{H}_k\) 在诸 \(\mathbf{H}_i\) 所成的集合中为极大，并证明 \(\mathbf{H}_k\) 是诸 \(\mathbf{H}_k\cap \mathbf{H}_i\) （其中 \(i\neq k\) ）的有限多个左陪集的并。）
\item
  给出一个交换环 A 及其理想 \(a, b_{1}, b_{2}, b_{3}\) 的例子，使得 \(a \notin b_{i} (i = 1, 2, 3)\) ，但 \(a = \bigcup_{i} b_{i}\) （利用 (a)）。
\end{enumerate}

\begin{enumerate}
\def\labelenumi{\arabic{enumi}.}
\setcounter{enumi}{1}
\item
  设 A 是一个不必交换的环，\(a, p_{1}, \ldots, p_{n}\) 是 A 的双边理想。假设 a 含于诸 \(p_{i}\) 的并，而且除至多两个以外所有 \(p_{i}\) 都是素理想（《代数》第 VIII 章 §8 习题 6）。证明 a 含于某个 \(p_{i}\) 之中。
\item
  在积环 \(\mathbf{A} = \mathbf{R}^{\mathbf{N}}\) 中，对每个 \(n \in \mathbf{N}\) ，设 \(\mathfrak{m}_n\) 是由那些 \(f\colon \mathbf{N} \to \mathbf{R}\) （满足 \(f(n) = 0\) ）组成的极大理想。证明 \(\mathfrak{a} = \mathbf{R}^{(\mathbf{N})}\) 是 \(\mathbf{A}\) 的一个理想，它含于诸 \(\mathfrak{m}_n\) 的并，但不含于这些极大理想中的任何一个。
\item
  在环 A 中，设 p 是素理想，a 是使 \(p \subset Aa\) 但 \(a \notin p\) 的元素。证明 p = ap。
\item
  \begin{enumerate}
  \def\labelenumii{(\alph{enumii})}
  \tightlist
  \item
    设 A 是 Noether 环，\(\mathfrak{a}\) 是 A 的理想。证明存在有限多个素理想 \(\mathfrak{p}_i\) ( \(1 \leqslant i \leqslant r\) )，使得 \(\mathfrak{p}_1\mathfrak{p}_2 \ldots \mathfrak{p}_r \subset \mathfrak{a}\) 。（用反证法，在包含 \(\mathfrak{a}\) 、异于 A 且不含任何素理想有限积的理想之中考虑一个极大元；
  \end{enumerate}
\end{enumerate}

另一方面注意，若理想 b 不是素的，则存在两个包含 b、异于 b 且满足 cd ⊂ b 的理想 c、d。）（参见第 IV 章。）

若 A 是整环且 \(a \neq 0\) ，则可设诸 \(p_{i} \neq 0\) 。

\begin{enumerate}
\def\labelenumi{(\alph{enumi})}
\setcounter{enumi}{1}
\tightlist
\item
  给出一个非 Noether 环 A 的例子，使 (a) 的结论不成立。（例如取 A 为 {[}0, 1{]} 上连续实值函数所成的环。）
\end{enumerate}

¶ 6. (a) 设 A 是一个环，a、b 是 A 的两个理想且 b 有限生成；证明若商环 A/a 与 A/b 都是 Noether 环，则 A/ab 也是 Noether 环（注意 b/ab 是有限生成的 (A/a)-模）。

\begin{enumerate}
\def\labelenumi{(\alph{enumi})}
\setcounter{enumi}{1}
\tightlist
\item
  证明若环 A 的每个素理想都是有限生成的，则 A 是 Noether 环。（用反证法，证明在 A 的非有限生成理想所成的集合中将有一个极大元 c，而由假设它不是素的；于是将有两个理想 \(a \supset c\) 、\(b \supset c\) ，它们异于 c 且满足 \(ab \subset c\) ；然后利用 (a)。）
\end{enumerate}

\begin{enumerate}
\def\labelenumi{\arabic{enumi}.}
\setcounter{enumi}{6}
\tightlist
\item
  设 \(\mathbf{A} = \prod_{i=1}^{n} \mathbf{A}_i\) 是一个有限环族的积，诸 \(\mathbf{A}_i\) 典范地等同于 \(\mathbf{A}\) 的理想。设 \(\mathbf{B}\) 是 \(\mathbf{A}\) 的一个子环，使得 \(\operatorname{pr}_i \mathbf{B} = \mathbf{A}_i\) （对 \(1 \leqslant i \leqslant n\) ）。
\end{enumerate}

\begin{enumerate}
\def\labelenumi{(\alph{enumi})}
\item
  证明若诸 \(\mathbf{A}_i\) 是 Noether 环（相应地，Artin 环），则 B 是 Noether 环（相应地，Artin 环）（参见《代数》第 VIII 章 § 2 习题 12）。
\item
  设 \(n_{i}\) 是 B 的理想，即 \(pr_{i}\) 在 B 上的限制的核。证明 B 的每个素理想 p 都含于某个 \(n_{i}\) 之中（利用命题 1）；推出对这个指标 i，\(pr_{i}\) \(p \neq A_{i}\) 。证明若每个 \(A_{i}\) 只含有限多个（\(k_{i}\) 个）极大理想，则 B 至多含 \(\sum_{i=1}^{n} k_{i}\) 个极大理想。
\item
  设 \(\Re(A)\) 、\(\Re(B)\) 分别是 \(A\) 与 \(B\) 的 Jacobson 根。证明 \(\Re(B) = B \cap \Re(A)\) 。若每个 \(A_i\) 只含有限多个极大理想，证明对每个整数 \(k > 1\) ，有 \(\mathrm{pr}_i((\Re(B))^k) = (\Re(A_i))^k\) （对所有 \(i\) ） ，并且 \((\Re(B))^k = (\Re(A))^k \cap B\) （把 \(\Re(B)\) 写成互异极大理想之积，并注意若 \(m\) 是 \(B\) 的极大理想且 \(\mathrm{pr}_i(m) = m_i\) 是 \(A_i\) 的极大理想，则 \(\mathrm{pr}_i(m^k) = m_i^k\) ）。
\end{enumerate}

\begin{enumerate}
\def\labelenumi{\arabic{enumi}.}
\setcounter{enumi}{7}
\tightlist
\item
  \begin{enumerate}
  \def\labelenumii{(\alph{enumii})}
  \tightlist
  \item
    设 \(\mathfrak{a}, \mathfrak{b}\) 是 \(\mathfrak{a}\) 环 \(\mathbf{A}\) 的两个互素理想。证明 \(\mathfrak{a}: \mathfrak{b} = \mathfrak{a}\) ，\(\mathfrak{b}: \mathfrak{a} = \mathfrak{b}\) 。若 \(\mathfrak{c}\) 是 \(\mathbf{A}\) 的理想且 \(\mathfrak{ac} \subset \mathfrak{b}\) ，则 \(\mathfrak{c} \subset \mathfrak{b}\) 。
  \end{enumerate}
\end{enumerate}

\begin{enumerate}
\def\labelenumi{(\alph{enumi})}
\setcounter{enumi}{1}
\item
  设 p、q 是两个互不包含的素理想；则 p: q = p 且 q: p = q。在多项式环 A = K{[}X, Y{]} （其中 K 是域）中给出两个主素理想 p、q 的例子，使二者互不包含且不互素。
\item
  设 \(a\) 是 A 中不是零因子的元素。证明若主理想 \(p = Aa\) 是素的，则关系 \(p = bc\) （关于 A 的两个理想 \(b\) 、\(c\) ）蕴含 \(b = A\) 或 \(c = A\) 。
\end{enumerate}

\begin{enumerate}
\def\labelenumi{\arabic{enumi}.}
\setcounter{enumi}{8}
\tightlist
\item
  \begin{enumerate}
  \def\labelenumii{(\alph{enumii})}
  \tightlist
  \item
    设 \((\mathfrak{a}_i)_{1 \leqslant i \leqslant n}\) 是环 A 的理想的一个有限族，这些理想两两互素，且没有哪个 \(\mathfrak{a}_i\) 形如 \(\mathfrak{bc} = \mathfrak{b} \cap \mathfrak{c}\) ，其中 \(\mathfrak{b}\) 与 \(\mathfrak{c}\) 是两个互素理想。若 \((\mathfrak{a}_j')_{1 \leqslant j \leqslant m}\) 是 A 的理想中另一个两两互素的族，使得 \(\mathfrak{a}_1\mathfrak{a}_2 \ldots \mathfrak{a}_n = \mathfrak{a}_1'\mathfrak{a}_2' \ldots \mathfrak{a}_m'\) ，证明 \(m \leqslant n\) ；若 \(m = n\) ，则存在一个置换 \(\pi\) （定义在 \([1, n]\) 上），使得 \(a_i' = a_{\pi(i)}\) （对 \(1 \leqslant i \leqslant n\) ） （利用命题 5 及《代数》第 VIII 章 § 1 习题 1 (d)）。
  \end{enumerate}
\end{enumerate}

\begin{enumerate}
\def\labelenumi{(\alph{enumi})}
\setcounter{enumi}{1}
\tightlist
\item
  设 A 是 Noether 环。证明 A 的每个理想 \(\mathfrak{a}\) 都等于有限多个两两互素理想的积，且其中没有一个是两个互素理想之积。（用反证法，在不具有此性质的理想所成的集合中考虑一个极大元。）
\end{enumerate}

\begin{enumerate}
\def\labelenumi{\arabic{enumi}.}
\setcounter{enumi}{9}
\tightlist
\item
  给出一个环 A 及其互异极大理想的无穷族 \((\mathfrak{m}_n)_{n\in \mathbf{N}}\) 的例子，使得 \(\bigcap_{n}\mathfrak{m}_{n} = 0\) ，但典范映射
\end{enumerate}

\[
\mathrm{A} \rightarrow \prod_ {n} (\mathrm{A} / \mathrm{m} _ {n})
\]

不是满射的。

\begin{enumerate}
\def\labelenumi{\arabic{enumi}.}
\setcounter{enumi}{10}
\tightlist
\item
  设 A 是一个不必交换的环，\(a, b\) 是 A 的两个双边理想且 \(a + b = A\) 。证明 \(a \cap b = ab + ba\) 。给出一个使 \(a \cap b \neq ab\) 的例子（考虑一个域上的 2 阶下三角矩阵环）。
\end{enumerate}

